\section[Notebook 1988-7: the idempotent of an index problem]{Notebook 1988-7: the idempotent of an index problem; two involutions\protect\verifie{Notebook1988_7}}\label{sec:quillen-1988-7}

\emph{Context.} On page~30 of \carnet{1988}{7} Quillen writes the idempotent representing the index of $p$ when $qp = 1 - x$ and $pq = 1 - y$ with $x$, $y$ small (of trace class), in a form adapted to the traces of odd powers; on page~33 he compares two involutions $F$ and $\varepsilon$.

\begin{encadre}
\begin{theoreme}[Quillen, notebook 1988-7, pp.~30, 33]\label{thm:quillen-1988-7}
In a ring, let $qp = 1 - x$ and $pq = 1 - y$. Then $px = yp$ and $xq = qy$. For $n \in \mathbf N$ put $\tilde q = (1 + x + \cdots + x^{2n-1})q$; then $\tilde qp = 1 - x^{2n}$, $p\tilde q = 1 - y^{2n}$, the row $(\tilde q, x^n)$ times the column $(p; x^n)$ is $1$, and
\[
e = \begin{pmatrix} p\\ x^n\end{pmatrix}\begin{pmatrix}\tilde q & x^n\end{pmatrix} = \begin{pmatrix} 1 - y^{2n} & px^n\\ x^n\tilde q & x^{2n}\end{pmatrix}
\]
is idempotent; $e - \operatorname{diag}(1, 0)$ has its entries in the ideal generated by $x$ and $y$. If $F^2 = \varepsilon^2 = 1$, then $(F - \varepsilon)^2 = 2 - F\varepsilon - \varepsilon F$ commutes with $F$ and with $\varepsilon$.
\end{theoreme}
\end{encadre}

\begin{proof}
$px = p - pqp = yp$. The geometric sum gives $\tilde qp = (\sum_{k<2n}x^k)(1 - x) = 1 - x^{2n}$, and $p\,x^k = y^kp$ gives $p\tilde q = (\sum_{k<2n} y^k)(1 - y) = 1 - y^{2n}$. A column times a row whose product the other way is $1$ is idempotent. The last identities are expansions.
\end{proof}

\emph{Verification in Lean.} The file \href{\leanblobmain/Notebook1988_7.lean}{\texttt{Notebook1988\_7.\allowbreak{}lean}} proves the theorem as \texttt{mul\_x}, \texttt{x\_mul}, \texttt{qt\_mul}, \texttt{mul\_qt}, \texttt{row\_mul\_col}, \texttt{isIdempotentElem\_e}, \texttt{e\_entries}, \texttt{sub\_sq\_eq}, \texttt{commute\_sub\_sq\_left} and \texttt{commute\_sub\_sq\_right}, in any ring. The trace formula for the index is not formalised. The axioms used are the three usual axioms of mathlib.

\litterature{Classical: the idempotent of the boundary map $K_1 \to K_0$ (Milnor, \emph{Introduction to algebraic $K$-theory}, \S 3) in the form used for Fredholm modules (Connes, 1985). The edition does not present it as new.}
